\documentclass[10pt,a4paper]{amsart}
\usepackage[margin=19mm]{geometry}
\usepackage{amsmath,amssymb,mathtools}
\usepackage[scr=rsfs,cal=euler]{mathalfa}
\usepackage{xfrac}
\usepackage[unicode,hidelinks,bookmarks=false]{hyperref}
\newcommand{\ie}{i.e.\ }
\newcommand{\cf}{cf.\ }
\newcommand{\resp}{respectively\ }
\newcommand{\Subsection}{\subsection}
\providecommand{\nobreakdash}{\nobreak\hbox{-}}
\setlength{\emergencystretch}{2em}
\multlinegap0pt
\usepackage{amssymb}
\usepackage{mathtools}
\usepackage[shortlabels]{enumitem}
\usepackage{algorithm}\usepackage{algpseudocode}
\newtheorem{assumption}{Assumption}
\newtheorem{lemma}{Lemma}
\newcommand{\RR}{\mathbb{R}}
\newcommand{\liph}{L_H}
\newcommand{\lmin}{\omega}
\newcommand{\norm}[1] {\left \| #1 \right \|}
\title{Instance-optimal stochastic convex optimization: Can we improve upon sample-average and robust stochastic approximation?}
\author{Liwei Jiang, Ashwin Pananjady}
\date{Source: arXiv:2603.25657v1 (CC BY 4.0)}
\begin{document}
\maketitle

\noindent\emph{Source and licence.} Instance-optimal stochastic convex optimization: Can we improve upon sample-average and robust stochastic approximation?. Version arXiv:2603.25657v1 (CC BY 4.0), distributed by arXiv under the Creative Commons Attribution 4.0 International licence, http://creativecommons.org/licenses/by/4.0/, as recorded in the arXiv licence metadata for this exact version and shown on the abstract page https://arxiv.org/abs/2603.25657v1. Full licence notice: \texttt{licenses/arxiv-2603.25657-CC-BY-4.0.txt}.\par\smallskip

\noindent\textit{Editorial scope.} Counted unit: the article's lemma lem:function\_justify\_liph (source0.tex lines 1737--1746). The article's complete original proof of it is delivered with it (source0.tex lines 1747--1793, one \texttt{proof} environment of 47 lines). No mathematics is written, repaired or re-derived: the statement and the proof are the article's own lines, in the article's own notation, and no retained line is substituted or re-flowed.  Retained uncounted as the source-local context the proof needs: lines 488--494.  Nothing was omitted from any retained span, so the package carries no omission record.  No retained line contains a citation, so the wrapper carries no bibliography and no bibliography entry.  No retained line contains a figure environment, an image-inclusion command or drawing code, so no image is delivered and the package declares no asset dependency.  Editorial formatting only, in addition: an AMS \texttt{amsart} preamble that restates the article's own declarations and macros used by the retained lines, namely the environments \texttt{assumption}, \texttt{lemma} on separate counters, together with the standard packages those lines need.  The retained environments print in this document as Assumption 1, Lemma 1 in order of appearance, whereas the article prints the same items with the numbering of its own full document.  The complete native source of the article as distributed in the arXiv e-print for this exact version is bundled unchanged as \texttt{sources/197232/source0.tex}.
\begin{assumption}\label{assum: non-quadratic}
	The population objective function $F$ is $\mu$-strongly convex and $L$-smooth on $\mathbb{R}^d$ with respect to the norm $\| \cdot \|$. We denote its minimizer by $x^\star$. Additionally, $F$ has $\liph$-Lipschitz Hessian in an instance-specific norm, meaning that for any $x, x' \in \RR^d$,
	\begin{align}\label{eqn:lip_hessian}
		\|\nabla^2 F(x^\star) ^{-1}(\nabla^2 F(x) - \nabla^2 F(x'))\| \le \liph \|x - x'\|.
	\end{align}
	We further define $\lmin := \inf_{\|v\| =1} \| \nabla^2 F(x^\star) v \|_*$, so  that $\|\nabla^2 F(x^\star) v\|_* \ge \lmin \|v\|$ for any $v \in \RR^d$.
\end{assumption}

\begin{lemma}\label{lem:function_justify_liph}
	For any $\liph' >0$, define a one-dimensional function $g$ by
	\begin{align*}
		g(t) = \begin{cases}
			\frac{1}{2} t^2 - \frac{\liph'^2}{6} t^4  + \frac{\liph'^4}{30} t^6 & |t| < \frac{1}{\liph'}\\
			\frac{8}{15\liph'}\left(|t| - \frac{1}{\liph'} \right) + \frac{9}{10\liph'^2} & |t| \ge \frac{1}{\liph'}
		\end{cases}
	\end{align*}
	and another function $G \colon \RR^d \rightarrow \RR$ by $G(x) = g(\|x\|_2).$  Then $ F(x) := G(x) + \frac{\mu}{2}\|x\|_2^2$ satisfies Assumption~\ref{assum: non-quadratic} with parameters $(\mu, 2+\mu, \frac{48\liph'}{15}).$ In particular, for any $\liph >0$, we can take $\liph' = \frac{ \liph}{4}$ and obtain a function satisfying Assumption~\ref{assum: non-quadratic} with parameters $(\mu, 2+\mu, \liph)$.
\end{lemma}

\begin{proof}
	It is straightforward to verify that $g$ is a convex function on $\RR$ and increasing on $[0,\infty)$. Therefore, $G$ is a convex function and $ F$ is $\mu$-strongly convex. Moreover, by construction, $g$ is $C^3$-smooth at the point $\frac{1}{\liph'}$, and as a result $G$ is $C^3$-smooth on $\mathbb{R}^d$. The Hessian of \(G\) is given by
	\begin{align*}
		\nabla^2 G(x)= u(\|x\|_2) I + v(\|x\|_2) xx^\top,
	\end{align*}		
	with
	\begin{align*}
		u(t)= \begin{cases}
			1-\frac{2\liph'^2}{3} t^2 + \frac{\liph'^4}{5}t^4 &\text{if } |t|<\frac{1}{\liph'}\\
			\frac{8}{15\liph' |t|} & 	\text{if } |t|\ge\frac{1}{\liph'}
		\end{cases}
	\end{align*}
	and 
	\begin{align*}
		v(t)= \begin{cases}
			-\frac{4}{3} \liph'^2 + \frac{4}{5} \liph'^4 t^2 &\text{if } |t|<\frac{1}{\liph'}\\
			- \frac{8}{15 \liph' t^3}  & 	\text{if }|t|\ge\frac{1}{\liph'}
		\end{cases}.
	\end{align*}
	By definition of $u$ and $v$, we have for any $x \in \RR^d$ that
	\begin{align*}
		\norm{\nabla^2 G(x)}_2 &\le |u(\|x\|_2)| + |v(\|x\|_2)| \|x\|_2^2 \le 2.
	\end{align*}
	By the mean value theorem, $\nabla G$ is $2$-Lipschitz. Since $\nabla  F(x) = \nabla G(x) + \mu x$, $\nabla F$ is $(2+ \mu)$-Lipschitz.
	
	Finally, we show that $\nabla^2  F$ is $\frac{48\liph'}{15}$-Lipschitz continuous. To this end, we denote the map $x \mapsto v(\|x\|_2) xx^\top$ by $h$. The Fr\'echet derivative of $h$ in the direction \(z\in\mathbb{R}^d\) is given by
	\begin{align*}
		Dh(x)[z] = w(\|x\|_2) \langle x, z \rangle xx^\top + v(\|x\|_2)\left(zx^\top + xz^\top\right),
	\end{align*}
	where 
	\begin{align*}
		w(t) = \begin{cases}
			\frac{8\liph'^4 }{5}  & \text{if } |t| < \frac{1}{\liph'}\\
			\frac{8}{5 \liph' t^5}  & \text{if } |t| \ge \frac{1}{\liph'}. 
		\end{cases}
	\end{align*}
	Using properties of the operator norm, we obtain the bound
	\begin{align*}
		\|Dh(x)[z]\|_2 & \le |w(\|x\|)| |\langle x, z \rangle| \|xx^\top\|_2 + |v(\|x\|)|\left(\|zx^\top\|_2+\|xz^\top\|_2\right)\\
		&\le \left(|w(\|x\|_2)| \|x\|_2^3 + 2|v(\|x\|_2)| \|x\|_2  \right)\|z\|_2.
	\end{align*}
	By definition of $w$ and $v$, it is straightforward to verify that for any $x \in \RR^d$,
	\begin{align*}
		|w(\|x\|_2)| \|x\|_2^3 + 2|v(\|x\|_2)| \|x\|_2 \le \frac{8 \liph'}{3}.
	\end{align*}
	Therefore, by mean value theorem, $h$ is $\frac{8 \liph'}{3}$-Lipschitz with respect to the operator norm. Additionally, one can verify that $|u'(t)| \le \frac{8 \liph'}{15}$  for all $t \in \RR$, so $x \mapsto u(\|x\|_2) I$ is $\frac{8 \liph'}{15}$ Lipschitz with respect to the operator norm. Combining the pieces, we see that $\nabla^2 G$ is $\frac{48\liph'}{15}$-Lipschitz with respect to the operator norm, and so is $\nabla^2  F$. The result follows.
\end{proof}

\end{document}
